\section{带符号循环图与通量坐标}\label{sec:preliminaries}

\subsection{带符号邻接矩阵与切换}

对整数 \(n\geq 8\)，令 \(G_n=C_n(1,2)\) 的顶点集为
\(\mathbb{Z}/n\mathbb{Z}\)，边为

\begin{equation}
\{i,i+1\}, \{i,i+2\}    (i \in \mathbb{Z}/n\mathbb{Z}).
\end{equation}

赋号是映射 \(\sigma:E(G_n)\to \{+1,-1\}\)。其带符号邻接矩阵
\(A_\sigma\) 是实对称矩阵：若 \(\{i,j\}\) 是边，则其 \(ij\) 元为该边符号，
否则为零。记

\begin{equation}
\rho(A_\sigma)=\max\{\lvert \lambda\rvert: \lambda \in \operatorname{spec}(A_\sigma)\}.
\end{equation}

对顶点符号函数 \(d:\mathbb{Z}/n\mathbb{Z}\to \{+1,-1\}\)，令
\(D=\operatorname{diag}(d_i)\)。把 \(\sigma_{ij}\) 换成
\(d_i\sigma_{ij}d_j\)，就把 \(A_\sigma\) 换成 \(DA_\sigma D\)，因而保持谱
不变。这个操作称为切换。

由于 \(G_n\) 连通且有 \(2n\) 条边，其圈空间维数为
\(2n-n+1=n+1\)。所以恰有 \(2^{n+1}\) 个切换类。我们采用适配于步长一
Hamilton 圈的圈通量坐标系。

\subsection{Hamilton 规范、三角形通量与四边形通量}

令 \(a_i\) 为步长一边 \(\{i,i+1\}\) 的符号，\(b_i\) 为步长二边
\(\{i,i+2\}\) 的符号。定义三角形通量和 Hamilton 和乐为

\begin{equation}
\begin{aligned}
\tau_i&=a_i a_{i+1} b_i, \\
\alpha&=\prod_{i=0}^{n-1} a_i.
\end{aligned}
\end{equation}

这 \(n\) 个三角形圈与步长一 Hamilton 圈构成一个圈基，因此
\((\tau_{0},\ldots,\tau_{n-1},\alpha)\) 决定切换类。

作切换，使所有步长一边均为正，仅允许从 \(n-1\) 到 \(0\) 跨越切口的边
例外。等价地，把向量延拓到整数格点，并施加扭曲边界条件

\begin{equation}
x_{i+n}=\alpha x_i.
\end{equation}

此时带符号算子具有局部形式

\begin{equation}
\label{eq:periodic-operator}
(A_\tau x)_i=x_{i-1}+x_{i+1}+\tau_{i-2}x_{i-2}+\tau_i x_{i+2}.
\end{equation}

特别选取的局部四边形通量字为

\begin{equation}
\label{eq:quadrilateral-flux}
Q_i=\tau_i \tau_{i+1}.
\end{equation}

对于 \(p\)-周期字 \(\tau\)，式~\eqref{eq:quadrilateral-flux} 蕴含
\(\prod_{i=0}^{p-1}Q_i=1\)。反过来，任何满足乘积为 \(+1\) 的字
\(Q\in\{+1,-1\}^p\) 恰有两个周期提升：选择 \(\tau_0=+1\) 或 \(-1\)，
再由 \(\tau_{i+1}=Q_i\tau_i\) 递推得到。

我们称

\begin{equation}
D(Q)=\{i:Q_i=+1\}
\end{equation}

为缺陷集，并记 \(d=\lvert D(Q)\rvert\)。还要使用循环局部统计量

\begin{equation}
\label{eq:local-defect-statistics}
\begin{aligned}
a&=\#\{i:Q_i=Q_{i+1}=+1\}, \\
b&=\#\{i:Q_i=Q_{i+2}=+1\}.
\end{aligned}
\end{equation}

当 \(n=8\) 时，图中另有两个步长二四边形，它们不属于式~
\eqref{eq:quadrilateral-flux} 编码的局部圈。这不会损失有限图信息：参数化每个
切换类的是圈基坐标 \((\tau,\alpha)\)，而不是单独的 \(Q\)；全赋号枚举使用
的正是这些坐标。

\subsection{Floquet 矩阵}\label{subsec:floquet-matrices}

设 \(\tau\) 的展示周期为 \(p\)。把整数指标写成 \(mp+r\)，其中
\(0\leq r<p\)，并寻找如下形式的 Bloch 向量：

\begin{equation}
x_{mp+r}=z^m v_r.
\end{equation}

代入式~\eqref{eq:periodic-operator} 得到 \(p\times p\) Laurent 矩阵
\(H_\tau(z)\)。更具体地，对从输出剩余类 \(r\) 到绝对源指标 \(j\) 的每个
跃迁，写 \(j=mp+s\)、\(0\leq s<p\)，并把跃迁系数乘以 \(z^m\) 加到
\((r,s)\) 元。短胞中不同跃迁发生重合时，这一约定会自动合并系数。

恒等式

\begin{equation}
\label{eq:fiber-transpose}
H_\tau(z)^T=H_\tau(z^{-1})
\end{equation}

来自对每个无向跃迁反向。因此，当 \(z\) 位于单位圆上时，\(H_\tau(z)\) 是
Hermitian 矩阵。

对无限周期相位，定义

\begin{equation}
\label{eq:squared-periodic-radius}
R(Q)=\sup_{\lvert z\rvert=1} \rho(H_\tau(z))^{2}.
\end{equation}

\(Q\) 的两个提升具有相同的值，所以这个记号没有歧义。注意，\(R(Q)\)
始终是谱半径的\textbf{平方}。

下文每个显式纤维证书都采用规范提升 \(\tau_0=+1\)、
\(\tau_{i+1}=Q_i\tau_i\)，以及有序纤维基
\((v_0,\ldots,v_{p-1})\)。另一个提升可由引理~
\ref{lem:operator-equivalences} 中的对角共轭和纤维重参数化得到，因此具有相同
的谱平方结论。

需要显式矩阵代表时，记
\(H_Q(z):=H_{\tau^{\mathrm{can}}}(z)\)，它表示由该规范提升和基构成的纤维。
记号 \(R(Q)\) 和矩 \(M_k(Q)\) 仍与提升无关。

若有限图的阶数为 \(n=pL\)，Hamilton 和乐为 \(\alpha\)，则其胞序列满足
\(u_{m+L}=\alpha u_m\)。酉胞平移的特征值恰好满足

\begin{equation}
\label{eq:finite-bloch-grid}
z^L=\alpha,
\end{equation}

有限矩阵分解为

\begin{equation}
\label{eq:finite-floquet-decomposition}
A_{pL,\alpha}  \cong  \bigoplus_{z^L=\alpha} H_\tau(z).
\end{equation}

式~\eqref{eq:finite-bloch-grid} 给出离散有限集，而式~
\eqref{eq:squared-periodic-radius} 中的 \(\lvert z\rvert=1\) 描述无限体积谱。
本文不会用其中一个替代另一个。

\subsection{算子等价}

下面记录分类周期相位时使用的等价关系。

\begin{lemma}[平移、反射与边符号取反]
\label{lem:operator-equivalences}
平移 \(\tau\)、按 \(\tau_i\mapsto\tau_{-i-2}\) 反射，或把 \(\tau\) 换成
\(-\tau\)，均保持 \(R(Q)\) 不变。在通量字上，反射为
\(Q_i\mapsto Q_{-i-3}\)。

\end{lemma}
\begin{proof}
令 \((T_rx)_i=x_{i-r}\)、\((Jx)_i=x_{-i}\)。直接代入式~
\eqref{eq:periodic-operator} 得

\begin{equation}
\begin{aligned}
T_r A_\tau T_r^{-1}&=A_{\text{平移后的 }\tau}, \\
J A_\tau J&=A_{\text{反射后的 }\tau}.
\end{aligned}
\end{equation}

二者都是酉共轭。对于全局边符号取反，令 \((Dx)_i=(-1)^ix_i\)。步长一边
的两个端点具有相反的 \(D\) 符号，而步长二边的端点符号相同。因此

\begin{equation}
A_{-\tau}=-D A_\tau D.
\end{equation}

谱被取负而其平方不变。在奇数展示胞中，纤维坐标从 \(z\) 变为 \(-z\)，
但完整单位圆保持不变。

\end{proof}
\begin{lemma}[能区折叠]
\label{lem:zone-folding}
若 \(\tau\) 的本原周期为 \(q\)，并展示在重复胞 \(p=mq\) 中，则

\begin{equation}
\label{eq:zone-folding}
H_p(z) \cong \bigoplus_{w^m=z} H_q(w).
\end{equation}

特别地，重复单位胞保持 \(R(Q)\) 不变。

\end{lemma}
\begin{proof}
把重复胞中的剩余类写成 \(r+kq\)，其中 \(0\leq r<q\)、\(0\leq k<m\)。
按酉内部平移 \(q\) 分解这个 \(mq\) 维纤维。在其 \(w\)-特征空间上，向量
具有形式

\begin{equation}
x_{r+kq}=w^k v_r.
\end{equation}

重复边界条件等价于 \(w^m=z\)。由于式~\eqref{eq:periodic-operator} 的所有
系数都是 \(q\)-周期的，每个跃迁都可提出 \(w^k\)，而在 \(v\) 上剩余作用
恰为 \(H_q(w)\)。\(w^m=z\) 的 \(m\) 个根给出两两正交的特征空间，其维数
之和为 \(mq\)。这就证明了式~\eqref{eq:zone-folding}，包括重数。

\end{proof}
\subsection{闭步矩}

对周期相位定义

\begin{equation}
\label{eq:moment-and-excess}
\begin{aligned}
M_k(Q)&=\operatorname{CT}_{z} \operatorname{tr}(H_Q(z)^{2k}), \\
F_k(Q)&=M_{k+1}(Q)-8M_k(Q).
\end{aligned}
\end{equation}

取常数项等于沿单位圆的归一化积分：

\begin{equation}
\label{eq:moment-integral}
M_k(Q)=\frac{1}{2\pi}\int_{0}^{2\pi}\sum_j\lambda_j(e^{i\theta})^{2k}\,d\theta.
\end{equation}

它也等于每个周期胞中长度为 \(2k\) 的带符号闭步之和。

\begin{lemma}[矩屏障]
\label{lem:moment-barrier}
若 \(R(Q)\leq 8\)，则对每个 \(k\geq 1\) 都有
\(M_{k+1}(Q)\leq 8M_k(Q)\)。等价地，

\begin{equation}
\label{eq:moment-barrier}
F_k(Q)>0  \Longrightarrow  R(Q)>8.
\end{equation}

\end{lemma}
\begin{proof}
在 \(R(Q)\leq 8\) 下，每个
\(y_j(\theta)=\lambda_j(e^{i\theta})^2\) 都位于 \([0,8]\)，所以
\(y_j^{k+1}\leq 8y_j^k\)。对 \(j\) 求和，再用式~
\eqref{eq:moment-integral} 积分即可。

\end{proof}
需要强调，式~\eqref{eq:moment-barrier} 是单向蕴含。一个非正超额量，甚至任意
有限列非正超额量，都不能证明 \(R(Q)\leq 8\)。

\subsection{特殊常数}

对偶数 \(n\geq 8\)，继承自文献 \cite{Suvagiya2026Signed} 的猜想阈值为

\begin{equation}
\label{eq:conjectured-threshold}
\begin{aligned}
\rho_-(n)&=2 \sqrt{\cos^{2}(\frac{\pi}{n})+\cos^{2}(\frac{2\pi}{n})}, \\
\rho_-(n)^{2}&=4+2\cos(\frac{2\pi}{n})+2\cos(\frac{4\pi}{n}).
\end{aligned}
\end{equation}

目标周期八谱边的平方为

\begin{equation}
\label{eq:target-edge}
\begin{aligned}
\eta&=4+\sqrt{10+2\sqrt{5}}, \\
\rho_*&=\sqrt{\eta}.
\end{aligned}
\end{equation}

这两个量作用不同：\(\rho_-(n)\) 是有限图的猜想下界，而 \(\eta\) 是反例
相位无限体积谱半径平方的精确值。
